\documentclass[12pt]{article}
\usepackage[utf8]{inputenc}
\usepackage[spanish]{babel}
\usepackage{amsmath}
\usepackage{amssymb}
\usepackage{geometry}
\usepackage{booktabs}
\geometry{margin=2.5cm}

\begin{document}
	
	\begin{center}
		{\Large \textbf{Guía 06 Ayudantía --- Ecuaciones}} \\[4pt]
		{\large Soluciones} \\[2pt]
		Universidad Santo Tomás --- Fundamentos de Matemática Básica
	\end{center}
	
	\bigskip
	
	%=============================================================
	\section*{Problema 1}
	%=============================================================
	
	\textbf{Ecuación:}
	\[
	0{,}85(12x + 150) - 60 = 900
	\]
	
	\textbf{Desarrollo:}
	\begin{align*}
		0{,}85(12x + 150) - 60 &= 900 \\
		10{,}2x + 127{,}5 - 60 &= 900 \\
		10{,}2x + 67{,}5       &= 900 \\
		10{,}2x                &= 832{,}5 \\
		x                      &= \frac{832{,}5}{10{,}2}
	\end{align*}
	\[
	\boxed{x \approx 81{,}62 \text{ hectáreas}}
	\]
	
	\textbf{Interpretación:} Se deben fertilizar aproximadamente $81{,}62$ hectáreas
	para que el cultivo reciba exactamente 900 kg de nitrógeno efectivo, descontadas
	las pérdidas por lixiviación del 15\%.
	
	\bigskip
	
	%=============================================================
	\section*{Problema 2}
	%=============================================================
	
	\textbf{Ecuación:}
	\[
	\frac{3}{4}(20x + 200) - (5x + 80) = 600
	\]
	
	\textbf{Desarrollo:}
	\begin{align*}
		\frac{3}{4}(20x + 200) - (5x + 80) &= 600 \\
		15x + 150 - 5x - 80                &= 600 \\
		10x + 70                           &= 600 \\
		10x                                &= 530
	\end{align*}
	\[
	\boxed{x = 53 \text{ horas de riego}}
	\]
	
	\textbf{Análisis de eficiencia:} El sistema requiere 53 horas de riego continuo,
	lo cual es elevado. Esto puede indicar baja eficiencia si se compara con métodos
	tecnificados (goteo o aspersión). Se recomienda revisar el caudal y la
	distribución del sistema.
	
	\bigskip
	
	%=============================================================
	\section*{Problema 3}
	%=============================================================
	
	\textbf{Ecuación:}
	\[
	\frac{5x + 40}{3} - \frac{2x - 10}{4} = 100
	\]
	
	\subsection*{a) Eliminación de denominadores}
	
	El $\text{mcm}(3,4) = 12$. Multiplicando ambos lados por $12$:
	\[
	4(5x + 40) - 3(2x - 10) = 1200
	\]
	
	\subsection*{b) Resolución}
	
	\begin{align*}
		20x + 160 - 6x + 30 &= 1200 \\
		14x + 190           &= 1200 \\
		14x                 &= 1010
	\end{align*}
	\[
	\boxed{x = \frac{1010}{14} \approx 72{,}14 \text{ mg/L}}
	\]
	
	\subsection*{c) Interpretación}
	
	La concentración base es aproximadamente $72{,}14$ mg/L. Este valor es la
	concentración mínima necesaria para que, al combinar las dos soluciones según
	el modelo, se alcance una concentración efectiva de 100 mg/L en el sistema.
	
	\bigskip
	
	%=============================================================
	\section*{Problema 4}
	%=============================================================
	
	\textbf{Ecuación:}
	\[
	15x + 200 = 8x + 500
	\]
	
	\subsection*{a) Punto de equilibrio}
	
	\begin{align*}
		15x - 8x &= 500 - 200 \\
		7x       &= 300
	\end{align*}
	\[
	\boxed{x = \frac{300}{7} \approx 42{,}86 \text{ toneladas}}
	\]
	
	\subsection*{b) Interpretación económica}
	
	Para menos de $42{,}86$ toneladas los costos superan los ingresos (pérdida);
	para más, los ingresos superan los costos (ganancia). Es el umbral mínimo de
	producción para que la actividad sea rentable.
	
	\bigskip
	
	%=============================================================
	\section*{Problema 5}
	%=============================================================
	
	\textbf{Ecuación:} $V = axt - bx + c$. \textbf{Despejar $x$:}
	
	\begin{align*}
		V - c     &= axt - bx \\
		V - c     &= x(at - b)
	\end{align*}
	\[
	\boxed{x = \frac{V - c}{at - b}, \quad at \neq b}
	\]
	
	\textbf{Interpretación:} $x$ es la cantidad necesaria para obtener un volumen
	efectivo $V$, dado un volumen inicial $c$, tasa de aplicación $a$, tiempo $t$
	y pérdidas por infiltración $b$.
	
	\bigskip
	
	%=============================================================
	\section*{Problema 6}
	%=============================================================
	
	\textbf{Ecuación:} $2(5x + 50) - 3x = 10x + 100$. \quad
	\textbf{Conclusión del estudiante:} $x = -10$.
	
	\subsection*{a) Verificación}
	
	Sustituimos $x = -10$:
	\begin{align*}
		\text{Lado izquierdo:} &\quad 2(5(-10)+50)-3(-10) = 2(0)+30 = 30 \\
		\text{Lado derecho:}   &\quad 10(-10)+100 = 0
	\end{align*}
	Como $30 \neq 0$, el resultado $x = -10$ \textbf{es incorrecto}.
	
	\subsection*{b) Resolución correcta}
	
	\begin{align*}
		10x + 100 - 3x &= 10x + 100 \\
		7x + 100       &= 10x + 100 \\
		-3x            &= 0
	\end{align*}
	\[
	\boxed{x = 0}
	\]
	
	\subsection*{c) Análisis de sentido físico}
	
	$x = 0$ indica que no se requiere insumo adicional. En agronomía puede
	interpretarse como que el campo ya posee la dosis óptima sin necesidad de
	intervención extra.
	
	\bigskip
	
	%=============================================================
	\section*{Problema 7}
	%=============================================================
	
	\textbf{Ecuación:}
	\[
	\frac{2(3x + 60)}{5} + \frac{x - 20}{2} = 150
	\]
	
	\subsection*{a) Resolución}
	
	El $\text{mcm}(5,2) = 10$. Multiplicando por $10$:
	\begin{align*}
		4(3x + 60) + 5(x - 20) &= 1500 \\
		12x + 240 + 5x - 100   &= 1500 \\
		17x + 140              &= 1500 \\
		17x                    &= 1360
	\end{align*}
	\[
	\boxed{x = 80}
	\]
	
	\subsection*{b) Interpretación como dosis de fertilizante}
	
	Una concentración base de $x = 80$ unidades permite que el modelo alcance
	la dosis total de 150 unidades requerida por el cultivo.
	
	\subsection*{c) Coherencia del modelo}
	
	El valor $x = 80 > 0$ es físicamente coherente y razonable dentro de los
	rangos típicos de dosificación agronómica.
	
	\bigskip
	
	%=============================================================
	\section*{Problema 8}
	%=============================================================
	
	\textbf{Ecuación:}
	\[
	0{,}9(10x + 300) - 0{,}2(5x + 100) = 1000
	\]
	
	\subsection*{a) Determinación de $x$}
	
	\begin{align*}
		9x + 270 - x - 20 &= 1000 \\
		8x + 250          &= 1000 \\
		8x                &= 750
	\end{align*}
	\[
	\boxed{x = \frac{750}{8} = 93{,}75}
	\]
	
	\subsection*{b) Interpretación en consumo energético}
	
	$x = 93{,}75$ representa las horas de operación o potencia requerida para
	entregar 1000 unidades de energía útil, con eficiencia del 90\% en la
	componente principal y pérdidas del 20\% en la secundaria.
	
	\bigskip
	
	%=============================================================
	\section*{Problema 9}
	%=============================================================
	
	\textbf{Ecuación:}
	\[
	\frac{3(4x + 120)}{2} - \frac{2(5x - 50)}{3} = 500
	\]
	
	\subsection*{a) Resolución}
	
	El $\text{mcm}(2,3) = 6$. Multiplicando por $6$:
	\begin{align*}
		9(4x + 120) - 4(5x - 50) &= 3000 \\
		36x + 1080 - 20x + 200   &= 3000 \\
		16x + 1280               &= 3000 \\
		16x                      &= 1720
	\end{align*}
	\[
	\boxed{x = 107{,}5}
	\]
	
	\subsection*{b) Interpretación en contexto productivo}
	
	$x = 107{,}5$ representa la cantidad de insumo que debe aplicarse para que
	el sistema productivo alcance las 500 unidades objetivo.
	
	\subsection*{c) Posibles errores de modelación}
	
	\begin{itemize}
		\item Los coeficientes deben calibrarse con datos reales del predio; errores
		en la estimación de pérdidas o rendimientos sesgan el resultado.
		\item Si $x$ representa una dosis, debe cotejarse con los rangos agronómicos
		recomendados para el cultivo específico.
		\item La ecuación asume linealidad, lo que puede ser una simplificación
		excesiva en sistemas con comportamiento no lineal.
	\end{itemize}
	
	\bigskip
	
	%=============================================================
	\section*{Problema 10}
	%=============================================================
	
	\textbf{Datos:} Sea $x$ el pago del día 1. Entonces:
	
	\begin{itemize}
		\item Día 1: $x$
		\item Día 2: $\dfrac{5}{4}x$
		\item Día 3: $\left(\dfrac{5}{4}\right)^{2}x = \dfrac{25}{16}x$
	\end{itemize}
	
	\textbf{Ecuación:}
	\[
	x + \frac{5}{4}x + \frac{25}{16}x = 219\,600
	\]
	
	Multiplicando por $16$:
	\begin{align*}
		16x + 20x + 25x &= 3\,513\,600 \\
		61x             &= 3\,513\,600
	\end{align*}
	\[
	\boxed{x = \$57\,600}
	\]
	
	\textbf{Pagos por día:}
	\[
	\text{Día 1: } \$57\,600 \qquad
	\text{Día 2: } \$72\,000 \qquad
	\text{Día 3: } \$90\,000
	\]
	
	\textbf{Verificación:} $57\,600 + 72\,000 + 90\,000 = \$219\,600$ \checkmark
	
	\bigskip
	
	%=============================================================
	\section*{Problema 11}
	%=============================================================
	
	\textbf{Datos:} muestra de $450$ plantas, $36$ infectadas; zona total: $200\,000$ plantas.
	
	\textbf{Proporción de infección:}
	\[
	p = \frac{36}{450} = 0{,}08 = 8\%
	\]
	
	\textbf{Estimación:}
	\[
	N_{\text{inf}} = 0{,}08 \times 200\,000
	\]
	\[
	\boxed{N_{\text{inf}} = 16\,000 \text{ plantas infectadas}}
	\]
	
	\textbf{Interpretación:} El 8\% del cultivo podría estar infectado, lo que
	amerita una intervención fitosanitaria inmediata para evitar la propagación.
	
	\bigskip
	
	%=============================================================
	\section*{Problema 12}
	%=============================================================
	
	\textbf{Datos:} Total $= 88$ profesionales; Grupo A $=$ Grupo B $+ 20$.
	
	Sea $B$ = profesionales en Grupo B:
	\begin{align*}
		(B + 20) + B &= 88 \\
		2B + 20      &= 88 \\
		2B           &= 68 \\
		B            &= 34
	\end{align*}
	\[
	\boxed{\text{Grupo A} = 34 + 20 = 54 \text{ profesionales}}
	\]
	
	\textbf{Verificación:} $54 + 34 = 88$ \checkmark
	
	\bigskip
	
	%=============================================================
	\section*{Resumen de Resultados}
	%=============================================================
	
	\begin{center}
		\renewcommand{\arraystretch}{1.8}
		\begin{tabular}{clc}
			\toprule
			\textbf{N\textsuperscript{o}} & \textbf{Resultado principal} & \textbf{Valor} \\
			\midrule
			1  & Hectáreas a fertilizar         & $x \approx 81{,}62$ ha \\
			2  & Horas de riego                 & $x = 53$ h \\
			3  & Concentración base             & $x \approx 72{,}14$ mg/L \\
			4  & Punto de equilibrio            & $x \approx 42{,}86$ ton \\
			5  & Despeje de $x$                 & $x = \dfrac{V-c}{at-b}$ \\[4pt]
			6  & Corrección del estudiante      & $x = 0$ \\
			7  & Dosis de fertilizante          & $x = 80$ \\
			8  & Consumo energético             & $x = 93{,}75$ \\
			9  & Insumo productivo              & $x = 107{,}5$ \\
			10 & Pago día 1                     & $\$57\,600$ \\
			11 & Plantas infectadas estimadas   & $16\,000$ plantas \\
			12 & Profesionales Grupo A          & $54$ profesionales \\
			\bottomrule
		\end{tabular}
	\end{center}
	
\end{document}